%% 
%% Copyright 2007-2025 Elsevier Ltd
%% 
%% This file is part of the 'Elsarticle Bundle'.
%% ---------------------------------------------
%% 
%% It may be distributed under the conditions of the LaTeX Project Public
%% License, either version 1.3 of this license or (at your option) any
%% later version.  The latest version of this license is in
%%    http://www.latex-project.org/lppl.txt
%% and version 1.3 or later is part of all distributions of LaTeX
%% version 1999/12/01 or later.
%% 
%% The list of all files belonging to the 'Elsarticle Bundle' is
%% given in the file `manifest.txt'.
%% 
%% Template article for Elsevier's document class `elsarticle'
%% with numbered style bibliographic references
%% SP 2008/03/01
%% $Id: elsarticle-template-num.tex 272 2025-01-09 17:36:26Z rishi $
%%
\documentclass[preprint,12pt]{elsarticle}

%% Use the option review to obtain double line spacing
%% \documentclass[authoryear,preprint,review,12pt]{elsarticle}

%% Use the options 1p,twocolumn; 3p; 3p,twocolumn; 5p; or 5p,twocolumn
%% for a journal layout:
%% \documentclass[final,1p,times]{elsarticle}
%% \documentclass[final,1p,times,twocolumn]{elsarticle}
%% \documentclass[final,3p,times]{elsarticle}
%% \documentclass[final,3p,times,twocolumn]{elsarticle}
%% \documentclass[final,5p,times]{elsarticle}
%% \documentclass[final,5p,times,twocolumn]{elsarticle}

%% For including figures, graphicx.sty has been loaded in
%% elsarticle.cls. If you prefer to use the old commands
%% please give \usepackage{epsfig}

%% The amssymb package provides various useful mathematical symbols
\usepackage{amssymb}
%% The amsmath package provides various useful equation environments.
\usepackage{amsmath}
%% The amsthm package provides extended theorem environments
%% \usepackage{amsthm}

%% The lineno packages adds line numbers. Start line numbering with
%% \begin{linenumbers}, end it with \end{linenumbers}. Or switch it on
%% for the whole article with \linenumbers.
%% \usepackage{lineno}

\journal{Nuclear Physics B}

\begin{document}

\begin{frontmatter}

%% Title, authors and addresses

%% use the tnoteref command within \title for footnotes;
%% use the tnotetext command for theassociated footnote;
%% use the fnref command within \author or \affiliation for footnotes;
%% use the fntext command for theassociated footnote;
%% use the corref command within \author for corresponding author footnotes;
%% use the cortext command for theassociated footnote;
%% use the ead command for the email address,
%% and the form \ead[url] for the home page:
%% \title{Title\tnoteref{label1}}
%% \tnotetext[label1]{}
%% \author{Name\corref{cor1}\fnref{label2}}
%% \ead{email address}
%% \ead[url]{home page}
%% \fntext[label2]{}
%% \cortext[cor1]{}
%% \affiliation{organization={},
%%             addressline={},
%%             city={},
%%             postcode={},
%%             state={},
%%             country={}}
%% \fntext[label3]{}

\title{}

%% use optional labels to link authors explicitly to addresses:
%% \author[label1,label2]{}
%% \affiliation[label1]{organization={},
%%             addressline={},
%%             city={},
%%             postcode={},
%%             state={},
%%             country={}}
%%
%% \affiliation[label2]{organization={},
%%             addressline={},
%%             city={},
%%             postcode={},
%%             state={},
%%             country={}}

\author{} %% Author name

%% Author affiliation
\affiliation{organization={},%Department and Organization
            addressline={}, 
            city={},
            postcode={}, 
            state={},
            country={}}

%% Abstract
\begin{abstract}
%% Text of abstract
Abstract text.
\end{abstract}

%%Graphical abstract
\begin{graphicalabstract}
%\includegraphics{grabs}
\end{graphicalabstract}

%%Research highlights
\begin{highlights}
\item Research highlight 1
\item Research highlight 2
\end{highlights}

%% Keywords
\begin{keyword}
%% keywords here, in the form: keyword \sep keyword

%% PACS codes here, in the form: \PACS code \sep code

%% MSC codes here, in the form: \MSC code \sep code
%% or \MSC[2008] code \sep code (2000 is the default)

\end{keyword}

\end{frontmatter}

%% Add \usepackage{lineno} before \begin{document} and uncomment 
%% following line to enable line numbers
%% \linenumbers

%% main text
%%

%% Use \section commands to start a section
\section{Example Section}
\label{sec1}
%% Labels are used to cross-reference an item using \ref command.

Section text. See Subsection \ref{subsec1}.

%% Use \subsection commands to start a subsection.
\subsection{Example Subsection}
\label{subsec1}

Subsection text.

%% Use \subsubsection, \paragraph, \subparagraph commands to 
%% start 3rd, 4th and 5th level sections.
%% Refer following link for more details.
%% https://en.wikibooks.org/wiki/LaTeX/Document_Structure#Sectioning_commands

\subsubsection{Mathematics}
%% Inline mathematics is tagged between $ symbols.
This is an example for the symbol $\alpha$ tagged as inline mathematics.

%% Displayed equations can be tagged using various environments. 
%% Single line equations can be tagged using the equation environment.
\begin{equation}
f(x) = (x+a)(x+b)
\end{equation}

%% Unnumbered equations are tagged using starred versions of the environment.
%% amsmath package needs to be loaded for the starred version of equation environment.
\begin{equation*}
f(x) = (x+a)(x+b)
\end{equation*}

%% align or eqnarray environments can be used for multi line equations.
%% & is used to mark alignment points in equations.
%% \\ is used to end a row in a multiline equation.
\begin{align}
 f(x) &= (x+a)(x+b) \\
      &= x^2 + (a+b)x + ab
\end{align}

\begin{eqnarray}
 f(x) &=& (x+a)(x+b) \nonumber\\ %% If equation numbering is not needed for a row use \nonumber.
      &=& x^2 + (a+b)x + ab
\end{eqnarray}

%% Unnumbered versions of align and eqnarray
\begin{align*}
 f(x) &= (x+a)(x+b) \\
      &= x^2 + (a+b)x + ab
\end{align*}

\begin{eqnarray*}
 f(x)&=& (x+a)(x+b) \\
     &=& x^2 + (a+b)x + ab
\end{eqnarray*}

%% Refer following link for more details.
%% https://en.wikibooks.org/wiki/LaTeX/Mathematics
%% https://en.wikibooks.org/wiki/LaTeX/Advanced_Mathematics

%% Use a table environment to create tables.
%% Refer following link for more details.
%% https://en.wikibooks.org/wiki/LaTeX/Tables
\begin{table}[t]%% placement specifier
%% Use tabular environment to tag the tabular data.
%% https://en.wikibooks.org/wiki/LaTeX/Tables#The_tabular_environment
\centering%% For centre alignment of tabular.
\begin{tabular}{l c r}%% Table column specifiers
%% Tabular cells are separated by &
  1 & 2 & 3 \\ %% A tabular row ends with \\
  4 & 5 & 6 \\
  7 & 8 & 9 \\
\end{tabular}
%% Use \caption command for table caption and label.
\caption{Table Caption}\label{fig1}
\end{table}


%% Use figure environment to create figures
%% Refer following link for more details.
%% https://en.wikibooks.org/wiki/LaTeX/Floats,_Figures_and_Captions
\begin{figure}[t]%% placement specifier
%% Use \includegraphics command to insert graphic files. Place graphics files in 
%% working directory.
\centering%% For centre alignment of image.
\includegraphics{example-image-a}
%% Use \caption command for figure caption and label.
\caption{Figure Caption}\label{fig1}
%% https://en.wikibooks.org/wiki/LaTeX/Importing_Graphics#Importing_external_graphics
\end{figure}


%% The Appendices part is started with the command \appendix;
%% appendix sections are then done as normal sections
\appendix
\section{Example Appendix Section}
\label{app1}

Appendix text.

%% For citations use: 
%%       \cite{<label>} ==> [1]

%%
Example citation, See \cite{lamport94}.

%% If you have bib database file and want bibtex to generate the
%% bibitems, please use
%%
%%  \bibliographystyle{elsarticle-num} 
%%  \bibliography{<your bibdatabase>}

%% else use the following coding to input the bibitems directly in the
%% TeX file.

%% Refer following link for more details about bibliography and citations.
%% https://en.wikibooks.org/wiki/LaTeX/Bibliography_Management

\begin{thebibliography}{00}

%% For numbered reference style
%% \bibitem{label}
%% Text of bibliographic item

@article{friedman2001greedy,
  title={{Greedy function approximation: a gradient boosting machine}},
  author={Friedman, Jerome H},
  journal={Annals of Statistics},
  pages={1189--1232},
volume={29},
number={5},
  year={2001},
  publisher={JSTOR}
}

@article{saha2023random,
  title={Random forests for spatially dependent data},
  author={Saha, Arkajyoti and Basu, Sumanta and Datta, Abhirup},
  journal={Journal of the American Statistical Association},
  volume={118},
  number={541},
  pages={665--683},
  year={2023},
  publisher={Taylor \& Francis}
}

@article{zhan2024neural,
  title={Neural networks for geospatial data},
  author={Zhan, Wentao and Datta, Abhirup},
  journal={Journal of the American Statistical Association},
  note={Published online: 24, June 2024},
  pages={1--21},
  year={2024},
  publisher={Taylor \& Francis}
}

@article{finley2024models,
  title={Models to support forest inventory and small area estimation using sparsely sampled {LiDAR}: A case study involving {G-LiHT} {LiDAR} in {T}anana, {A}laska},
  author={Finley, Andrew O and Andersen, Hans-Erik and Babcock, Chad and Cook, Bruce D and Morton, Douglas C and Banerjee, Sudipto},
  journal={Journal of Agricultural, Biological and Environmental Statistics},
  pages={1--28},
  volume={29},
  year={2024},
  publisher={Springer}
}

@article{osgood2023statistical,
  title={A Statistical Review of Template Model Builder: A Flexible Tool for Spatial Modelling},
  author={Osgood-Zimmerman, Aaron and Wakefield, Jon},
  journal={International Statistical Review},
  volume={91},
  OPTnumber={2},
  pages={318--342},
  year={2023},
  publisher={Wiley Online Library}
}

@article{banerjee2024finite,
  title={Finite population survey sampling: An unapologetic Bayesian perspective},
  author={Banerjee, Sudipto},
  journal={Sankhya A},
  pages={1--30},
volume={86},
  year={2024},
  publisher={Springer}
}


@article{chan2020bayesian,
  title={Bayesian inference for finite populations under spatial process settings},
  author={Chan-Golston, Alec M and Banerjee, Sudipto and Handcock, Mark S},
  journal={Environmetrics},
  volume={31},
  number={3},
  pages={e2606},
  year={2020},
  publisher={Wiley Online Library}
}

@article{tan2019bayesian,
  title={Bayesian additive regression trees and the general {BART} model},
  author={Tan, Yaoyuan Vincent and Roy, Jason},
  journal={Statistics in Medicine},
  volume={38},
  number={25},
  pages={5048--5069},
  year={2019},
  publisher={Wiley Online Library}
}

@article{Dorie, 
   url={https://cran.r-project.org/web/packages/dbarts/vignettes/gibbs_sampler_mixture_model.pdf}, 
   title={Building a Gibbs sampler with dbarts},
   journal={\texttt{dbarts}}, 
   author={Dorie, Vincent},
   year = {2020}} 

@article{polson2013bayesian,
  title={Bayesian inference for logistic models using P{\'o}lya--Gamma latent variables},
  author={Polson, Nicholas G and Scott, James G and Windle, Jesse},
  journal={Journal of the American Statistical Association},
  volume={108},
  number={504},
  pages={1339--1349},
  year={2013},
  publisher={Taylor \& Francis}
}

@article{utazi2018high,
  title={High resolution age-structured mapping of childhood vaccination coverage in low and middle income countries},
  author={Utazi, C Edson and Thorley, Julia and Alegana, Victor A and Ferrari, Matthew J and Takahashi, Saki and Metcalf, C Jessica E and Lessler, Justin and Tatem, Andrew J},
  journal={Vaccine},
  volume={36},
  number={12},
  pages={1583--1591},
  year={2018},
  publisher={Elsevier}
}

@article{burstein2019mapping,
  title={Mapping 123 million neonatal, infant and child deaths between 2000 and 2017},
  author={Burstein, Roy and Henry, Nathaniel J and Collison, Michael L and Marczak, Laurie B and Sligar, Amber and Watson, Stefanie and Marquez, Neal and Abbasalizad-Farhangi, Mahdieh and Abbasi, Masoumeh and Abd-Allah, Foad and others},
  journal={Nature},
  volume={574},
  number={7778},
  pages={353--358},
  year={2019},
  publisher={Nature Publishing Group UK London}
}



@article{mei2007standard,
  title={{Standard deviation of anthropometric Z-scores as a data quality assessment tool using the 2006 WHO growth standards: a cross country analysis}},
  author={Mei, Zuguo and Grummer-Strawn, Laurence M},
  journal={Bulletin of the World Health Organization},
  volume={85},
  pages={441--448},
  year={2007},
  publisher={SciELO Public Health}
}

@article{kassie2019exploring,
  title={Exploring the association of anthropometric indicators for under-five children in {E}thiopia},
  author={Kassie, Gashu Workneh and Workie, Demeke Lakew},
  journal={BMC Public Health},
  volume={19},
  number={1},
  pages={1--6},
  year={2019},
  publisher={BioMed Central}
}

@article{dezeure2015high,
  title={High-dimensional inference: confidence intervals, p-values and {R}-software hdi},
  author={Dezeure, Ruben and B{\"u}hlmann, Peter and Meier, Lukas and Meinshausen, Nicolai},
  journal={Statistical Science},
  pages={533--558},
volume={30},
number={4},
  year={2015},
  publisher={JSTOR}
}

@book{rao2015small,
  title={Small Area Estimation},
  author={Rao, John NK and Molina, Isabel},
  year={2015},
  publisher={John Wiley \& Sons}
}

@article{battese1988error,
  title={An error-components model for prediction of county crop areas using survey and satellite data},
  author={Battese, George E and Harter, Rachel M and Fuller, Wayne A},
  journal={Journal of the American Statistical Association},
  volume={83},
  number={401},
  pages={28--36},
  year={1988},
  publisher={Taylor \& Francis}
}

@article{fay1979estimates,
  title={Estimates of income for small places: an application of {James-Stein} procedures to census data},
  author={Fay, Robert E and Herriot, Roger A},
  journal={Journal of the American Statistical Association},
  volume={74},
  number={366a},
  pages={269--277},
  year={1979},
  publisher={Taylor \& Francis}
}

@article{knorr2002block,
  title={On block updating in {M}arkov random field models for disease mapping},
  author={Knorr-Held, Leonhard and Rue, H{\aa}vard},
  journal={Scandinavian Journal of Statistics},
  volume={29},
  number={4},
  pages={597--614},
  year={2002},
  publisher={Wiley Online Library}
}

@article{martins2013bayesian,
  title={Bayesian computing with {INLA}: new features},
  author={Martins, Thiago G and Simpson, Daniel and Lindgren, Finn and Rue, H{\aa}vard},
  journal={Computational Statistics and Data Analysis},
  volume={67},
  pages={68--83},
  year={2013},
  publisher={Elsevier}
}

@article{sparapani2021nonparametric,
  title={{Nonparametric machine learning and efficient computation with Bayesian additive regression trees: the BART R package}},
  author={Sparapani, Rodney and Spanbauer, Charles and McCulloch, Robert},
  journal={Journal of Statistical Software},
  volume={97},
  pages={1--66},
  year={2021}
}

@article{eddelbuettel2011rcpp,
  title={{Rcpp: Seamless R and C++ Integration}},
  author={Eddelbuettel, Dirk and Fran{\c{c}}ois, Romain},
  journal={Journal of Statistical Software},
  volume={40},
  pages={1--18},
  year={2011}
}

@book{eddelbuettel2013seamless,
  title={Seamless {R} and {C}++ Integration with {R}cpp},
  author={Eddelbuettel, Dirk},
  year={2013},
  publisher={Springer}
}

@article{eddelbuettel2018extending,
  title={{Extending R with C++: a brief introduction to Rcpp}},
  author={Eddelbuettel, Dirk and Balamuta, James Joseph},
  journal={The American Statistician},
  volume={72},
  number={1},
  pages={28--36},
  year={2018},
  publisher={Taylor \& Francis}
}




@book{dellaportas2003introduction,
  title={An introduction to {MCMC}},
  author={Dellaportas, Petros and Roberts, Gareth O},
  journal={Spatial Statistics and Computational Methods},
  pages={1--41},
  year={2003},
  publisher={Springer}
}

@book{diggle2019model,
  title={Model-Based Geostatistics for Global Public Health: Methods and Applications},
  author={Diggle, Peter J and Giorgi, Emanuele},
  year={2019},
  publisher={CRC Press}
}

@article{lindstrom2014flexible,
  title={A flexible spatio-temporal model for air pollution with spatial and spatio-temporal covariates},
  author={Lindstr{\"o}m, Johan and Szpiro, Adam A and Sampson, Paul D and Oron, Assaf P and Richards, Mark and Larson, Tim V and Sheppard, Lianne},
  journal={Environmental and Ecological Statistics},
  volume={21},
  pages={411--433},
  year={2014},
  publisher={Springer}
}

@article{ZERAATPISHEH2022105723,
title = {Improving the spatial prediction of soil organic carbon using environmental covariates selection: A comparison of a group of environmental covariates},
journal = {CATENA},
volume = {208},
pages = {105723},
year = {2022},
issn = {0341-8162},
doi = {https://doi.org/10.1016/j.catena.2021.105723},
url = {https://www.sciencedirect.com/science/article/pii/S0341816221005816},
author = {Mojtaba Zeraatpisheh and Younes Garosi and Hamid {Reza Owliaie} and Shamsollah Ayoubi and Ruhollah Taghizadeh-Mehrjardi and Thomas Scholten and Ming Xu},
keywords = {Soil organic carbon, Spatial prediction, Environmental covariates, Time-series vegetation indices, Machine learning, Uncertainty},
abstract = {}
}

@article{chipman1998bayesian,
  title={{Bayesian CART model search}},
  author={Chipman, Hugh A and George, Edward I and McCulloch, Robert E},
  journal={Journal of the American Statistical Association},
  volume={93},
  number={443},
  pages={935--948},
  year={1998},
  publisher={Taylor \& Francis}
}

@article{gneiting2007strictly,
  title={Strictly proper scoring rules, prediction, and estimation},
  author={Gneiting, Tilmann and Raftery, Adrian E},
  journal={Journal of the American Statistical Association},
  volume={102},
  number={477},
  pages={359--378},
  year={2007},
  publisher={Taylor \& Francis}
}

@article{breiman2001random,
  title={Random forests},
  author={Breiman, Leo},
  journal={Machine Learning},
  volume={45},
  number={1},
  pages={5--32},
  year={2001},
  publisher={Springer}
}

@article{georganos2021geographical,
  title={Geographical random forests: a spatial extension of the random forest algorithm to address spatial heterogeneity in remote sensing and population modelling},
  author={Georganos, Stefanos and Grippa, Tais and Niang Gadiaga, Assane and Linard, Catherine and Lennert, Moritz and Vanhuysse, Sabine and Mboga, Nicholus and Wolff, El{\'e}onore and Kalogirou, Stamatis},
  journal={Geocarto International},
  volume={36},
  number={2},
  pages={121--136},
  year={2021},
  publisher={Taylor \& Francis}
}

@article{daw2023reds,
  title={{REDS: Random ensemble deep spatial prediction}},
  author={Daw, Ranadeep and Wikle, Christopher K},
  journal={Environmetrics},
  volume={34},
  number={1},
  pages={e2780},
  year={2023},
  publisher={Wiley Online Library}
}



@article{goldstein2015peeking,
  title={{Peeking inside the black box: Visualizing statistical learning with plots of individual conditional expectation}},
  author={Goldstein, Alex and Kapelner, Adam and Bleich, Justin and Pitkin, Emil},
  journal={Journal of Computational and Graphical Statistics},
  volume={24},
  number={1},
  pages={44--65},
  year={2015},
  publisher={Taylor \& Francis}
}



@article{ren2018,
title = {{Maternal exposure to ambient PM10 during pregnancy increases the risk of congenital heart defects: Evidence from machine learning models}},
journal = {Science of The Total Environment},
volume = {630},
pages = {1-10},
year = {2018},
issn = {0048-9697},
doi = {https://doi.org/10.1016/j.scitotenv.2018.02.181},
url = {https://www.sciencedirect.com/science/article/pii/S0048969718305746},
author = {Zhoupeng Ren and Jun Zhu and Yanfang Gao and Qian Yin and Maogui Hu and Li Dai and Changfei Deng and Lin Yi and Kui Deng and Yanping Wang and Xiaohong Li and Jinfeng Wang},
keywords = {Congenital heart defects, Machine learning, Air pollution, Birth defects, Particulate matter},
abstract = {Previous research suggested an association between maternal exposure to ambient air pollutants and risk of congenital heart defects (CHDs), though the effects of particulate matter ≤10μm in aerodynamic diameter (PM10) on CHDs are inconsistent. We used two machine learning models (i.e., random forest (RF) and gradient boosting (GB)) to investigate the non-linear effects of PM10 exposure during the critical time window, weeks 3–8 in pregnancy, on risk of CHDs. From 2009 through 2012, we carried out a population-based birth cohort study on 39,053 live-born infants in Beijing. RF and GB models were used to calculate odds ratios for CHDs associated with increase in PM10 exposure, adjusting for maternal and perinatal characteristics. Maternal exposure to PM10 was identified as the primary risk factor for CHDs in all machine learning models. We observed a clear non-linear effect of maternal exposure to PM10 on CHDs risk. Compared to 40μgm−3, the following odds ratios resulted: 1) 92μgm−3 [RF: 1.16 (95% CI: 1.06, 1.28); GB: 1.26 (95% CI: 1.17, 1.35)]; 2) 111μgm−3 [RF: 1.04 (95% CI: 0.96, 1.14); GB: 1.04 (95% CI: 0.99, 1.08)]; 3) 124μgm−3 [RF: 1.01 (95% CI: 0.94, 1.10); GB: 0.98 (95% CI: 0.93, 1.02)]; 4) 190μgm−3 [RF: 1.29 (95% CI: 1.14, 1.44); GB: 1.71 (95% CI: 1.04, 2.17)]. Overall, both machine models showed an association between maternal exposure to ambient PM10 and CHDs in Beijing, highlighting the need for non-linear methods to investigate dose-response relationships.}
}

@article{macharia2019sub,
  title={{Sub national variation and inequalities in under-five mortality in Kenya since 1965}},
  author={Macharia, Peter M and Giorgi, Emanuele and Thuranira, Pamela N and Joseph, Noel K and Sartorius, Benn and Snow, Robert W and Okiro, Emelda A},
  journal={BMC Public Health},
  volume={19},
  number={1},
  pages={1--12},
  year={2019},
  publisher={BioMed Central}
}

@article{lindgren2011explicit,
  title={{An explicit link between Gaussian fields and Gaussian Markov random fields: the stochastic partial differential equation approach}},
  author={Lindgren, Finn and Rue, H{\aa}vard and Lindstr{\"o}m, Johan},
  journal={Journal of the Royal Statistical Society: Series B},
  volume={73},
  number={4},
  pages={423--498},
  year={2011},
  publisher={Wiley Online Library}
}

@article{wakefield2019estimating,
  title={{Estimating under-five mortality in space and time in a developing world context}},
  author={Wakefield, Jon and Fuglstad, Geir-Arne and Riebler, Andrea and Godwin, Jessica and Wilson, Katie and Clark, Samuel J},
  journal={Statistical Methods in Medical Research},
  volume={28},
  number={9},
  pages={2614--2634},
  year={2019},
  publisher={SAGE Publications Sage UK: London, England}
}

@article{rue2017bayesian,
  title={{Bayesian computing with INLA: a review}},
  author={Rue, H{\aa}vard and Riebler, Andrea and S{\o}rbye, Sigrunn H and Illian, Janine B and Simpson, Daniel P and Lindgren, Finn K},
  journal={Annual Review of Statistics and Its Application},
  volume={4},
  pages={395--421},
  year={2017},
  publisher={Annual Reviews}
}

@article{spanbauer2021nonparametric,
  title={{Nonparametric machine learning for precision medicine with longitudinal clinical trials and Bayesian additive regression trees with mixed models}},
  author={Spanbauer, Charles and Sparapani, Rodney},
  journal={Statistics in Medicine},
  volume={40},
  number={11},
  pages={2665--2691},
  year={2021},
  publisher={Wiley Online Library}
}

@article{Krueger2020,
title = {{A new spatial count data model with Bayesian additive regression trees for accident hot spot identification}},
journal = {Accident Analysis and Prevention},
volume = {144},
pages = {105623},
year = {2020},
issn = {0001-4575},
doi = {https://doi.org/10.1016/j.aap.2020.105623},
url = {https://www.sciencedirect.com/science/article/pii/S0001457520306680},
author = {Rico Krueger and Prateek Bansal and Prasad Buddhavarapu},
keywords = {Accident analysis, Site ranking, Spatial count data modelling, Negative binomial model, Bayesian additive regression trees, Pólya-Gamma data augmentation},
abstract = {The identification of accident hot spots is a central task of road safety management. Bayesian count data models have emerged as the workhorse method for producing probabilistic rankings of hazardous sites in road networks. Typically, these methods assume simple linear link function specifications, which, however, limit the predictive power of a model. Furthermore, extensive specification searches are precluded by complex model structures arising from the need to account for unobserved heterogeneity and spatial correlations. Modern machine learning (ML) methods offer ways to automate the specification of the link function. However, these methods do not capture estimation uncertainty, and it is also difficult to incorporate spatial correlations. In light of these gaps in the literature, this paper proposes a new spatial negative binomial model which uses Bayesian additive regression trees to endogenously select the specification of the link function. Posterior inference in the proposed model is made feasible with the help of the Pólya-Gamma data augmentation technique. We test the performance of this new model on a crash count data set from a metropolitan highway network. The empirical results show that the proposed model performs at least as well as a baseline spatial count data model with random parameters in terms of goodness of fit and site ranking ability.}
}

@article{muller2007spatially,
  title={{A spatially-adjusted Bayesian additive regression tree model to merge two datasets}},
  author={M{\"u}ller, Peter and Shih, Ya-Chen Tina and Zhang, Song},
  journal={Bayesian Analysis},
  volume={2},
  number={3},
  pages={611--633},
  year={2007},
  publisher={International Society for Bayesian Analysis}
}

@article{lesage07,
title = {{A matrix exponential spatial specification}},
journal = {Journal of Econometrics},
volume = {140},
number = {1},
pages = {190-214},
year = {2007},
note = {Analysis of spatially dependent data},
issn = {0304-4076},
doi = {https://doi.org/10.1016/j.jeconom.2006.09.007},
url = {https://www.sciencedirect.com/science/article/pii/S0304407606002284},
author = {James P. LeSage and R. {Kelley Pace}},
keywords = {Spatial autoregression, Bayesian, Maximum likelihood, Log-determinants, Matrix exponentials, Model comparison},
abstract = {We introduce the matrix exponential as a way of modelling spatially dependent data. The matrix exponential spatial specification (MESS) simplifies the log-likelihood allowing a closed form solution to the problem of maximum-likelihood estimation, and greatly simplifies the Bayesian estimation of the model. The MESS can produce estimates and inferences similar to those from conventional spatial autoregressive models, but has analytical, computational, and interpretive advantages. We present maximum likelihood and Bayesian approaches to the estimation of this spatial model specification along with methods of model comparisons over different explanatory variables and spatial specifications.}
}

@article{shi2021digital,
  title={{Digital mapping of zinc in urban topsoil using multisource geospatial data and random forest}},
  author={Shi, Tiezhu and Hu, Xianjun and Guo, Long and Su, Fenzheng and Tu, Wei and Hu, Zhongwen and Liu, Huizeng and Yang, Chao and Wang, Jingzhe and Zhang, Jie and others},
  journal={Science of the Total Environment},
  volume={792},
  pages={148455},
  year={2021},
  publisher={Elsevier}
}

@article{davies2016optimal,
  title={{Optimal spatial prediction using ensemble machine learning}},
  author={Davies, Molly Margaret and Van Der Laan, Mark J},
  journal={The International Journal of Biostatistics},
  volume={12},
  number={1},
  pages={179--201},
  year={2016},
  publisher={De Gruyter}
}

@article{osgood2018mapping,
  title={{Mapping child growth failure in Africa between 2000 and 2015}},
  author={Osgood-Zimmerman, Aaron and Millear, Anoushka I and Stubbs, Rebecca W and Shields, Chloe and Pickering, Brandon V and Earl, Lucas and Graetz, Nicholas and Kinyoki, Damaris K and Ray, Sarah E and Bhatt, Samir and others},
  journal={Nature},
  volume={555},
  number={7694},
  pages={41--47},
  year={2018},
  publisher={Nature Publishing Group}
}

@article{Fuglstad19,
author = {Geir-Arne Fuglstad and Daniel Simpson and Finn Lindgren and HÃ¥vard Rue},
title = {Constructing Priors that Penalize the Complexity of {G}aussian Random Fields},
journal = {Journal of the American Statistical Association},
volume = {114},
number = {525},
pages = {445-452},
year  = {2019},
publisher = {Taylor & Francis},
doi = {10.1080/01621459.2017.1415907},

URL = { 
        https://doi.org/10.1080/01621459.2017.1415907
    
},
eprint = { 
        https://doi.org/10.1080/01621459.2017.1415907
    
}

}

@article{uwiringiyimana2022bayesian,
  title={{Bayesian geostatistical modelling of stunting in Rwanda: risk factors and spatially explicit residual stunting burden}},
  author={Uwiringiyimana, Vestine and Osei, Frank and Amer, Sherif and Veldkamp, Antonie},
  journal={BMC Public Health},
  volume={22},
  number={1},
  pages={1--14},
  year={2022},
  publisher={BioMed Central}
}

@article{hubin2016estimating,
  title={{Estimating the marginal likelihood with Integrated nested Laplace approximation (INLA)}},
  author={Hubin, Aliaksandr and Storvik, Geir},
  journal={arXiv preprint arXiv:1611.01450},
  year={2016}
}

@article{bhatt2017improved,
  title={{Improved prediction accuracy for disease risk mapping using Gaussian process stacked generalization}},
  author={Bhatt, Samir and Cameron, Ewan and Flaxman, Seth R and Weiss, Daniel J and Smith, David L and Gething, Peter W},
  journal={Journal of The Royal Society Interface},
  volume={14},
  number={134},
  pages={20170520},
  year={2017},
  publisher={The Royal Society}
}



@article{rue2009approximate,
  title={{Approximate Bayesian inference for latent Gaussian models by using integrated nested Laplace approximations}},
  author={Rue, H{\aa}vard and Martino, Sara and Chopin, Nicolas},
  journal={Journal of the Royal Statistical Society: Series B},
  volume={71},
  number={2},
  pages={319--392},
  year={2009},
  publisher={Wiley Online Library}
}

@article{chipman2010bart,
  title={{BART: Bayesian additive regression trees}},
  author={Chipman, Hugh A and George, Edward I and McCulloch, Robert E},
  journal={The Annals of Applied Statistics},
  volume={4},
  number={1},
  pages={266--298},
  year={2010},
  publisher={Institute of Mathematical Statistics}
}

@article{ribeiro2006analysis,
  title={Analysis of geostatistical data},
  author={Ribeiro Jr, Paulo J and Diggle, Peter J},
  journal={The geoR package, version},
  pages={1--6},
  year={2006}
}

@article{marsan2008,
  title={Extending earthquakes' reach through cascading},
  author={Marsan, David and Lengline, Olivier},
  journal={Science},
  volume={319},
  number={5866},
  pages={1076--1079},
  year={2008},
  publisher={American Association for the Advancement of Science}
}

@inproceedings{zhou2013,
  title={Learning triggering kernels for multi-dimensional hawkes processes},
  author={Zhou, Ke and Zha, Hongyuan and Song, Le},
  booktitle={International conference on machine learning},
  pages={1301--1309},
  year={2013},
  organization={PMLR}
}



@article{laub2015hawkes,
  title={Hawkes processes},
  author={Laub, Patrick J and Taimre, Thomas and Pollett, Philip K},
  journal={arXiv preprint arXiv:1507.02822},
  year={2015}
}

@article{gomez2018markov,
  title={Markov chain {M}onte {C}arlo with the integrated nested {L}aplace approximation},
  author={G{\'o}mez-Rubio, Virgilio and Rue, H{\aa}vard},
  journal={Statistics and Computing},
  volume={28},
  number={5},
  pages={1033--1051},
  year={2018},
  publisher={Springer}
}


@article{Chiang2022,
title = {Hawkes process modeling of COVID-19 with mobility leading indicators and spatial covariates},
journal = {International Journal of Forecasting},
volume = {38},
number = {2},
pages = {505-520},
year = {2022},
issn = {0169-2070},
doi = {https://doi.org/10.1016/j.ijforecast.2021.07.001},
url = {https://www.sciencedirect.com/science/article/pii/S0169207021001126},
author = {Wen-Hao Chiang and Xueying Liu and George Mohler},
keywords = {COVID-19 forecasting, Hawkes processes, Mobility indices, Spatial covariate, Demographic covariate, Epidemic modeling},
abstract = {Hawkes processes are used in statistical modeling for event clustering and causal inference, while they also can be viewed as stochastic versions of popular compartmental models used in epidemiology. Here we show how to develop accurate models of COVID-19 transmission using Hawkes processes with spatial-temporal covariates. We model the conditional intensity of new COVID-19 cases and deaths in the U.S. at the county level, estimating the dynamic reproduction number of the virus within an EM algorithm through a regression on Google mobility indices and demographic covariates in the maximization step. We validate the approach on both short-term and long-term forecasting tasks, showing that the Hawkes process outperforms several models currently used to track the pandemic, including an ensemble approach and an SEIR-variant. We also investigate which covariates and mobility indices are most important for building forecasts of COVID-19 in the U.S.}
}

@article{dempster1977maximum,
  title={Maximum likelihood from incomplete data via the EM algorithm},
  author={Dempster, Arthur P and Laird, Nan M and Rubin, Donald B},
  journal={Journal of the Royal Statistical Society: Series B (Methodological)},
  volume={39},
  number={1},
  pages={1--22},
  year={1977},
  publisher={Wiley Online Library}
}

@article{logothetis1999,
  author={Logothetis, A. and Krishnamurthy, V.},
  journal={IEEE Transactions on Signal Processing}, 
  title={Expectation maximization algorithms for MAP estimation of jump Markov linear systems}, 
  year={1999},
  volume={47},
  number={8},
  pages={2139-2156},
  doi={10.1109/78.774753}}
  
@book{Bishop2006,
author = {Bishop, Christopher M.},
title = {Pattern Recognition and Machine Learning (Information Science and Statistics)},
year = {2006},
isbn = {0387310738},
publisher = {Springer-Verlag},
address = {Berlin, Heidelberg}
}  

@article{Robbins51,
author = {Herbert Robbins and Sutton Monro},
title = {{A Stochastic Approximation Method}},
volume = {22},
journal = {The Annals of Mathematical Statistics},
number = {3},
publisher = {Institute of Mathematical Statistics},
pages = {400 -- 407},
year = {1951},
doi = {10.1214/aoms/1177729586},
URL = {https://doi.org/10.1214/aoms/1177729586}
}


@article{Hoffman13,
author = {Hoffman, Matthew D. and Blei, David M. and Wang, Chong and Paisley, John},
title = {Stochastic Variational Inference},
year = {2013},
issue_date = {January 2013},
publisher = {JMLR.org},
volume = {14},
number = {1},
issn = {1532-4435},
abstract = {We develop stochastic variational inference, a scalable algorithm for approximating posterior distributions. We develop this technique for a large class of probabilistic models and we demonstrate it with two probabilistic topic models, latent Dirichlet allocation and the hierarchical Dirichlet process topic model. Using stochastic variational inference, we analyze several large collections of documents: 300K articles from Nature, 1.8M articles from The New York Times, and 3.8M articles from Wikipedia. Stochastic inference can easily handle data sets of this size and outperforms traditional variational inference, which can only handle a smaller subset. (We also show that the Bayesian nonparametric topic model outperforms its parametric counterpart.) Stochastic variational inference lets us apply complex Bayesian models to massive data sets.},
journal = {J. Mach. Learn. Res.},
month = {may},
pages = {1303–1347},
numpages = {45},
keywords = {variational inference, stochastic optimization, Bayesian nonparametrics, Bayesian inference, topic models}
}

@article{tierney1986accurate,
  title={Accurate approximations for posterior moments and marginal densities},
  author={Tierney, Luke and Kadane, Joseph B},
  journal={Journal of the American Statistical Association},
  volume={81},
  number={393},
  pages={82--86},
  year={1986},
  publisher={Taylor \& Francis}
}

@article{Cappe09,
author = {Cappé, Olivier and Moulines, Eric},
title = {On-line expectation–maximization algorithm for latent data models},
journal = {Journal of the Royal Statistical Society: Series B (Statistical Methodology)},
volume = {71},
number = {3},
pages = {593-613},
keywords = {Adaptive algorithms, Expectation–maximization, Latent data models, Mixture of regressions, On-line estimation, Polyak–Ruppert averaging, Stochastic approximation},
doi = {https://doi.org/10.1111/j.1467-9868.2009.00698.x},
url = {https://rss.onlinelibrary.wiley.com/doi/abs/10.1111/j.1467-9868.2009.00698.x},
eprint = {https://rss.onlinelibrary.wiley.com/doi/pdf/10.1111/j.1467-9868.2009.00698.x},
abstract = {Summary.  We propose a generic on-line (also sometimes called adaptive or recursive) version of the expectation–maximization (EM) algorithm applicable to latent variable models of independent observations. Compared with the algorithm of Titterington, this approach is more directly connected to the usual EM algorithm and does not rely on integration with respect to the complete-data distribution. The resulting algorithm is usually simpler and is shown to achieve convergence to the stationary points of the Kullback–Leibler divergence between the marginal distribution of the observation and the model distribution at the optimal rate, i.e. that of the maximum likelihood estimator. In addition, the approach proposed is also suitable for conditional (or regression) models, as illustrated in the case of the mixture of linear regressions model.},
year = {2009}
}




\end{thebibliography}
\end{document}

\endinput
%%
%% End of file `elsarticle-template-num.tex'.
